% Audit chapter: the paper section is input, then thesis-only implementation notes, validation and the confirmatory results (generated tables).
\section{Counterfactual lesion audit}
\label{sec:audit}
% PLAN (about 1.5 pages). The definition of A is already in methods.tex (Lesion-absorption audit); this section states the protocol and its limits.
%   5.1 Definition and reading of A: slope of the change in the bone estimate on the inserted signal; A=0 keeps the lesion, A=1 removes it; A=0 is
%       reached by doing nothing, so A is always reported next to a suppression gain; A can leave [0,1].
%   5.2 Sites and lesions: lesions from held-out CTs (library built from the TEST split, s4/lesion_lib.py), six sites per image at alpha=1, sites drawn
%       once with a fixed seed and shared by all models (s4/audit_real.py), lung masks from a published network, bone overlap from independent masks.
%   5.3 Why the audit is not circular: the regulariser pastes lesions to teach invariance and the audit pastes lesions to measure it. Mitigations:
%       held-out lesion library, different image domain, independent rib masks, and evidence that does not use insertion (H4' in SAP_v2.md).
%   5.4 Variants: exact-truth audit on the Qure-projection test split (secondary), SZCH audit (development evidence only), amplitude sweep.
%   5.5 Statistics (SAP_v2.md H1', H3'): image-cluster bootstrap, slope of A on log2(1+w), Holm over contrasts.
\subsection{Definition}
For a suppressor $f$ that outputs the bone image, an image $x$ and a lesion signal $n\ge0$ in composite-image units, the absorption at amplitude
$\alpha$ is
\begin{equation}
A=\frac{\big\langle f(x+\alpha n)-f(x),\;\alpha n\big\rangle}{\big\langle \alpha n,\;\alpha n\big\rangle},
\label{eq:absorption}
\end{equation}
with both inner products taken over the lesion support. $A=0$ means that the lesion stays in the soft-tissue image and $A=1$ that it was removed with
the bone; $A$ is a regression slope and can fall outside $[0,1]$. A model that does nothing has $A=0$, so $A$ is reported together with a suppression
gain and never alone. Two suppression measures are paired with $A$: the bone-region PSNR on the held-out SZCH pairs and, on VinDr-RibCXR, the
reduction of gradient energy on rib-boundary pixels (the rib benchmark).

\subsection{Sites and lesions}
Lesions come from the CTs of the LIDC-IDRI test split only, so no audited lesion was seen in training. On each real radiograph six sites are drawn at random
positions inside the lung fields at $\alpha=1$; the sites depend only on the image, its lung mask, the lesion library and a fixed seed, and are shared by
every model. Lung masks come from the PSPNet anatomical segmentation model of TorchXRayVision \citep{cohen2022torchxrayvision}, eroded by
49~px at $1024$~px so that a site lies inside the lung field; the model comes from the ChestX-Det project \citep{lian2021chestxdet}, whose paper does not fully
document its training data. The library holds 296 nodules
(4.0 to 37.6~mm, median 7.1~mm); six sites per image give 1,470 sites on VinDr-RibCXR, 1,482 on JSRT and 450 on SZCH. Because the PSPNet masks
differ slightly between GPUs, evaluations run on cloud GPUs read the lung masks exported on the machine of the final table, so every model is
scored on identical sites.
The bone overlap of a site is measured with independent information, never with a model output: the radiologist rib masks of VinDr-RibCXR
\citep{nguyen2021ribcxr} for the confirmatory audit, and the real dual-energy bone signal on SZCH. A site \emph{overlaps bone} when at least 15\,\% of its
support lies on an annotated rib. On JSRT no bone reference exists; the plan named the consensus of the bone maps of ours and of the reference
model, but the audit took the overlap from the reference model's bone map alone (secondary analysis; recorded as a deviation).

\subsection{Why the audit is not circular}
The twin regulariser pastes lesions to teach invariance, and the audit pastes lesions to measure it, so a model could pass the audit by learning the
insertion artefact. Four design choices limit this. The audit lesions come from CTs that were never used for training; the audit images are real
radiographs, whereas the training composites are synthetic; the confirmatory audit uses rib masks that no model has seen; and the claim that the
regulariser helps is also tested with evidence that does not use insertion: contrast of real nodules on JSRT and Node21 and detection recall of real
lesions grouped by bone overlap (Section~\ref{sec:experiments}). These choices limit the circularity but cannot remove it: a model that recognises
pasted signal in general could still score well, which is why the evidence without insertion is reported beside the audit (Section~\ref{sec:limitations}).

\subsection{Variants}
\emph{Exact-truth audit.} On the Qure-projection test split the lesion is added to the soft component, the true bone image is known, and the
bone error inside the lesion support can be measured against it in addition to $A$. \emph{SZCH audit.} It ran on all 75 SZCH images during development
and is reported as development evidence only (Section~\ref{sec:data}). \emph{Amplitude sweep.} $A$ is recomputed for
$\alpha\in\{0.25,0.5,1,2,4\}$ on the 450 SZCH sites (development data, supplementary material). For $\alpha\le1$ the absorption is flat
(v1 model 0.301--0.306, no-twin model 0.215--0.218, reference model 0.134--0.136), so the leak is linear in the lesion amplitude; at $\alpha=2$ and
4 it falls for every model, where the added signal starts to saturate at the top of the intensity range.

\subsection{Statistics}
Site-level differences between models are paired; confidence intervals are image-cluster bootstrap intervals; the dose-response over the twin weight is
the slope of $A$ on $\log_2(1+w)$ with an image-cluster bootstrap interval, and the contrasts of ours against the reference model and against the
no-twin control are adjusted with the Holm procedure (SAP\_v2.md). Bootstrap intervals use \num{10000} resamples for mean absorption and contrasts,
\num{2000} for the overlap interaction and 500 for the dose-response slope, always resampling whole images with replacement.


\section{Implementation notes}
The audit is implemented in \path{s4/audit_real.py}. For each radiograph it draws six sites with a random generator seeded once (seed 7):
for every site a lesion is drawn from the library and a centre is drawn among the pixels of the eroded lung mask; the lesion patch, rendered
at $512$~px from a LIDC-IDRI test CT, is scaled to the $1024$~px frame and added in image units. The sites depend only on the image, its lung
mask, the library and the seed, so every model, run separately, is scored on the same sites, and the paired differences between models are
site by site. A site whose support (pixels above 2\,\% of the lesion maximum) has fewer than 30 pixels is skipped. For each model the bone
estimate of the radiograph is computed once, and for each site the estimate of the radiograph with the lesion; $A$ is the regression slope of
Equation~\ref{eq:absorption} over the support, and a leak term records the change of the bone estimate far from the lesion. The bone overlap of
a site is the fraction of its support that lies on bone: the radiologist rib masks on VinDr-RibCXR, the strongest 20\,\% of the dual-energy
bone signal on SZCH-X-Rays, and a thresholded bone map of the reference model on JSRT and BDCXR-3257.

Two practical points matter for reproducibility. First, the lung masks come from a neural network whose output differs slightly between GPUs,
which moves the sites; evaluations that ran on cloud GPUs therefore read the lung masks exported on the machine that produced the final
tables, and the identity of the sites was checked file by file (Chapter~\ref{ch:ethics}). Second, the audit of a three-model ensemble runs
the three members on each modified image, so it is the most expensive evaluation of the study; on one T4 GPU a batch of four models, two of them
ensembles, took 16 minutes for the 245 VinDr-RibCXR images.

\section{Validation of the audit}
\paragraph{Controls.} The identity has $A=0$ by construction, at every site and on every data set, which the tables confirm. It also removes
no bone, which is why absorption is always shown next to a suppression measure.

\paragraph{Exact truth.} On the Qure-projection test split the lesion is added to the soft component and the true bone image is known
(Table~\ref{tab:app_qp}). The ranking of the models is the same as on real radiographs: ours absorbs 0.102, the matched control 0.220 and the
reference model 0.199, and the contrast of the lesion in our soft estimate is 0.99 of its true contrast against 0.94 for the control. The cost
of the regulariser shows here as well: the bone error of ours (0.0041) is higher than that of the control (0.0029), though lower than that of
the reference model on these images (0.0082).

\paragraph{Linearity.} On the development SZCH sites the amplitude of the inserted lesion was varied over $\alpha\in\{0.25,0.5,1,2,4\}$
(Table~\ref{tab:app_amplitude}). Up to $\alpha=1$ the absorption is flat for every model (v1 0.301--0.306, no twin 0.215--0.218, reference
model 0.134--0.136), so $A$ measures a property of the model, not of the lesion amplitude; above that it falls for all of them, where the
added signal begins to saturate at the top of the intensity range. The audit is run at $\alpha=1$.

\section{Confirmatory audit on VinDr-RibCXR}
On the 1,470 sites of VinDr-RibCXR (Tables~\ref{tab:audit_ribcxr} and~\ref{tab:audit_contrasts}) ours absorbs $A=0.051$ $[0.045,0.057]$,
the reference model 0.149 $[0.138,0.160]$ and the no-twin control 0.178 $[0.167,0.188]$. The paired differences are $-0.098$
$[-0.108,-0.089]$ and $-0.127$ $[-0.134,-0.119]$, Holm-adjusted $p<0.001$ for both: H3$'$ passes. The interaction between sites on an annotated
rib and the other sites is $0.003$ $[-0.017,0.027]$ against the reference model, so the advantage is the same with and without rib overlap. On
JSRT the same pattern holds ($A$ 0.048, 0.118 and 0.121) with a larger advantage at bone-overlapping sites ($-0.060$ $[-0.074,-0.046]$), but
there the overlap was defined by the reference model's own bone map. The SZCH audit, development evidence, agrees (0.054, 0.130 and 0.195).
Figure~\ref{fig:audit_maps} shows one site on a held-out SZCH image, chosen by a fixed rule.
\begin{table}[htbp]
\centering\small
\caption{Counterfactual lesion audit on VinDr-RibCXR: mean absorption per method with image-cluster bootstrap intervals (lower is better).}
\label{tab:audit_ribcxr}
\fittable{tables/audit_ribcxr.tex}
\end{table}
\begin{table}[htbp]
\centering\small
\caption{Paired site-level contrasts on the three audit sets, Holm over the two contrasts, and the interaction with bone overlap
($\geq$15\,\% vs.\ $<$15\,\%). SZCH is development evidence; the JSRT overlap comes from the reference model's bone map.}
\label{tab:audit_contrasts}
\fittable{tables/audit_contrasts.tex}
\end{table}
\begin{figure}[htbp]
\centering
\safefig{figures/audit_maps.pdf}
\caption{The audit on one held-out SZCH-X-Rays image (CC BY-NC 4.0): a 7\,mm nodule from a LIDC test CT is added; the change of each bone
estimate (red = lesion signal moved into the bone image) and the resulting soft-tissue image. Eight placements were drawn with a fixed seed and
the one with the median difference between ours and the reference model is shown.}
\label{fig:audit_maps}
\end{figure}

\section{Failure cases and limits of the audit}
The audit measures the leak of an added signal under a linear, additive model of image formation. Real lesions change attenuation inside the
patient, interact with scatter and with the tone curve of the detector, and may differ in texture from the CT patches used here; $A$ is
therefore a property of the suppressor for additive signals, not a measure of clinical detectability. Lesions near the lung edge can extend
beyond the eroded lung mask, very small lesions have supports near the 30-pixel limit where the regression slope is noisy, and saturated pixels
clip the added signal, which is why the amplitude is kept at $\alpha=1$. The bone overlap of a site is only as good as its reference: rib masks
on VinDr-RibCXR, the dual-energy signal on SZCH, and a model's bone map elsewhere. Finally, the audit and the twin regulariser both insert
lesions; the audit uses lesions from CTs never seen in training and real radiographs never used for training, and Chapter~\ref{ch:benchmark}
reports evidence that uses no insertion at all, but a model that learned to recognise pasted signal in general could still score well.
